\chapter{Аффинный KO6-бимодуль и канонический опорный угол}

Проверка координатной алгебры показала, что полная матричная алгебра
$M_{300}(\mathbb C)$ не может служить координатной. Поэтому вместо нового
перебора перечитаем три уже доказанных результата проекта:
\begin{enumerate}
 \item четырёхсостоянийное аффинное меню имеет каноническое разложение
 $\mathbb C^4=\mathbf1\oplus\mathbf3$;
 \item семейная цепь имеет явную координатную алгебру
 $\mathbb R_0\oplus M_3(\mathbb R)_G\oplus\mathbb C_2$ и точное
 18-мерное KO6-завершение;
 \item наблюдаемый пакет без правого нейтрино имеет размерность 15.
\end{enumerate}
Их совместное использование приводит ровно к размерности 300 и устраняет
искусственное прямое сложение отдельного опорного блока.

\section{Каноническое разложение четырёх состояний}

Пусть $U_g$ --- перестановочное представление всех 24 элементов
$S_4\simeq\operatorname{AGL}(2,2)$. Тогда
\begin{equation}
 P_1=\frac1{24}\sum_{g\in S_4}U_g=\frac14J_4,
 \qquad
 P_3=I_4-P_1.
 \label{eq:v5-affine-p1-p3}
\end{equation}
Проекторы имеют ранги один и три и коммутируют со всем аффинным действием.
Матрица
\begin{equation}
 V=
 \begin{pmatrix}
  1/\sqrt2&-1/\sqrt2&0&0\\
  1/\sqrt6&1/\sqrt6&-2/\sqrt6&0\\
  1/\sqrt{12}&1/\sqrt{12}&1/\sqrt{12}&-3/\sqrt{12}
 \end{pmatrix}
 \label{eq:v5-affine-coisometry}
\end{equation}
удовлетворяет
\begin{equation}
 VV^T=I_3,
 \qquad
 V^TV=P_3,
 \qquad
 V(1,1,1,1)^T=0.
\end{equation}
Для каждой перестановки существует ортогональная матрица
\begin{equation}
 R_g=VU_gV^T,
 \qquad
 R_gV=VU_g.
\end{equation}
Следовательно, $V$ является не выбранным базисом трёх поколений, а
сплетающим оператором канонического аффинного триплета. Его остаточная
$O(3)$-свобода совпадает с семейной калибровочной свободой среднего узла.

\section{Единый бимодуль размерности 300}

В старой семейной геометрии заменим левую кратность три на полное
четырёхсостоянийное меню:
\begin{equation}
 (0,0)^{\oplus4}
 \xrightarrow{\ X\ }
 (G,0)
 \xrightarrow{\ \Phi I_3\ }
 (G,2).
 \label{eq:v5-affine-family-chain}
\end{equation}
Частичные размерности узлов равны
\begin{equation}
 4+3+3=10.
\end{equation}
После умножения на 15-мерный наблюдаемый пакет и KO6-удвоения получаем
\begin{equation}
 \boxed{
 2(4+3+3)\cdot15=300.
 }
 \label{eq:v5-affine-m300-dimension}
\end{equation}
Это в точности прежняя амальгамированная размерность, но теперь 30-мерный
опорный сектор не добавлен извне: он является равномерным синглетом первой
четырёхмерной вершины и его $J$-образом.

Семейной координатной алгеброй остаётся
\begin{equation}
 \mathcal A_{\mathrm{fam}}
 =\mathbb R_0\oplus M_3(\mathbb R)_G\oplus\mathbb C_2.
\end{equation}
На частичной половине она действует блоками
\begin{equation}
 a_0I_4,
 \qquad A_G,
 \qquad A_G,
\end{equation}
а на сопряжённой половине --- транспонированными бимодульными метками.

\section{KO6, высота и ориентированный дифференциал}

Положим
\begin{equation}
 X=\rho V,
 \qquad Y=\Phi I_3,
\end{equation}
и определим частичный оператор Дирака
\begin{equation}
 D_p=
 \begin{pmatrix}
 0&X^T&0\\
 X&0&\overline\Phi I_3\\
 0&\Phi I_3&0
 \end{pmatrix}.
\end{equation}
Градуировка и высота имеют вид
\begin{equation}
 \gamma_p=\operatorname{diag}(I_4,-I_3,I_3),
 \qquad
 h_p=\operatorname{diag}(-P_3,0_3,I_3).
 \label{eq:v5-affine-height}
\end{equation}
Сопряжённая половина несёт противоположные градуировку и высоту, а $J$
обменивает половины с комплексным сопряжением.

Прямая проверка полного базиса алгебры даёт нулевые дефекты для
\begin{align}
 D&=D^\dagger,&
 \{D,\gamma\}&=0,&
 JD&=DJ,&
 J\gamma&=-\gamma J,\\
 [\pi(a),\pi^o(b)]&=0,&
 [[D,\pi(a)],\pi^o(b)]&=0.
\end{align}
Ориентированный дифференциал удовлетворяет
\begin{equation}
 D=d+d^\dagger,
 \qquad
 [h,d]=d,
 \qquad
 JdJ^{-1}=d^\dagger.
\end{equation}
Таким образом, забытая конструкция является настоящим 20-мерным
семейным KO6-бимодулем, а 300 возникает только после физической кратности
15.

\section{Опорный сектор как канонический угол}

На полном 20-мерном семейном модуле определим
\begin{equation}
 P_{\mathrm{ref}}
 =\operatorname{diag}(P_1,0_3,0_3;
                       P_1,0_3,0_3),
 \qquad
 P_{\mathrm{phys}}=I-P_{\mathrm{ref}}.
 \label{eq:v5-affine-reference-projector}
\end{equation}
В вакууме $X=V$ проектор $P_{\mathrm{ref}}$ коммутирует с алгеброй,
противоположной алгеброй, $J$, $\gamma$, $h$ и $D$. Следовательно, это
редуцирующий угол, а не вручную выбранные 15 координатных состояний.

После 15-кратного наблюдаемого умножения ранги равны
\begin{equation}
 \rank P_{\mathrm{ref}}=30,
 \qquad
 \rank P_{\mathrm{phys}}=270.
\end{equation}
Полный базовый оператор имеет восьмимерное ядро. Ограничение на физический
угол имеет шестимерное полное KO6-ядро, то есть после умножения на 15
\begin{equation}
 \dim\ker D_{\mathrm{phys}}=90.
\end{equation}
Физическая половина поэтому содержит ровно
\begin{equation}
 \boxed{45}
\end{equation}
лёгких частичных состояний. Нормированный условный след физического угла
следует из одного родительского следа с фиксированным весом
\begin{equation}
 \frac{270}{300}=\frac9{10}.
\end{equation}

\begin{theorem}[Канонический физический угол]
Аффинный проектор $P_3$ превращает амальгамированное разложение
$30+270$ в редукцию одного KO6-бимодуля. Физический угол инвариантен при
$J$, градуировке, координатной алгебре и вакуумном операторе Дирака и
содержит ровно 45 лёгких частичных состояний.
\end{theorem}

Это каноническая редукция, но не динамическая масса опорного синглета.
Родительская теория по-прежнему должна объяснить, почему наблюдаемые
фермионы читаются в углу $P_{\mathrm{phys}}$, а не во всём носителе.

\section{Кривизна и устойчивость смешивания с опорой}

Новая высота даёт частичные блоки кривизны
\begin{equation}
 \mathcal F_H=
 \operatorname{diag}\left(
 P_3-X^TX,
 XX^T-|\Phi|^2I_3,
 (|\Phi|^2-1)I_3
 \right).
 \label{eq:v5-affine-curvature-blocks}
\end{equation}
Нулевая кривизна достигается при
\begin{equation}
 X=V,
 \qquad |\Phi|=1.
\end{equation}
На двенадцати вещественных компонентах прямоугольной матрицы $X$ и двух
компонентах $\Phi$ сокращённый гессиан имеет точный спектр
\begin{equation}
 \left\{
 0^{(4)},
 \left(\frac43\right)^{(3)},
 \left(\frac{16}3\right)^{(5)},
 \frac{32-8\sqrt7}{3},
 \frac{32+8\sqrt7}{3}
 \right\}.
 \label{eq:v5-affine-hessian}
\end{equation}
Три новых собственных значения $4/3$ принадлежат как раз вариациям
$XP_1$, смешивающим опорный синглет со средней вершиной. Следовательно,
канонический физический угол не только инвариантен в вакууме: отклонения от
него штрафуются тем же квадратом кривизны. Сигнатура равна
\begin{equation}
 (10_+,4_0,0_-).
\end{equation}

\section{Исправление нормировки полного следа}

На полном 300-мерном носителе множитель наблюдаемого пакета сокращается, и
точно
\begin{equation}
 \tau_{300}(\mathcal F_H^2)
 =\frac1{10}\Tr_{10}(\mathcal F_{H,p}^2).
 \label{eq:v5-affine-parent-trace}
\end{equation}
Сокращённый потенциал предыдущей главы был записан как
\begin{equation}
 V_{\mathrm{red}}
 =\frac13\Tr_{10}(\mathcal F_{H,p}^2),
\end{equation}
поэтому
\begin{equation}
 \boxed{
 \tau_{300}(\mathcal F_H^2)=\frac3{10}V_{\mathrm{red}}.
 }
 \label{eq:v5-affine-trace-scale}
\end{equation}
Для нулевого множества и сигнатуры гессиана это лишь общий положительный
масштаб. Но единичный дефектный импульс должен быть повторно нормирован из
того же полного пространственно-конечного действия. Поэтому прежние
точные радиусы $5/6$, $2/3$ и связанная с ними масса пока сохраняются
только в сокращённой нормировке, а не как завершённое предсказание
$\tau_{300}$.

\section{Оставшаяся граница}

Проект восстановил семейную координатную алгебру, KO6-бимодуль, опорный
угол и устойчивость его смешиваний. Однако 15-мерный пакет пока входит как
кратность и представление калибровочной группы. Он ещё не объединён с
$\mathcal A_{\mathrm{fam}}$ в одну ассоциативную координатную алгебру.

Наивное тензорное произведение алгебр запрещено: матричные блоки типа
$M_3\otimes M_3$ порождают увеличенные унитарные группы вместо требуемого
произведения семейной и наблюдаемой симметрий. Следующий гейт должен
классифицировать амальгамы или бимодульные размещения, сохраняющие старую
алгебру Стандартной модели без такого увеличения.

Археология стандарт-модельной ветви подсказывает более узкий первый
кандидат. В обычной конечной спектральной тройке поколения входят как
кратности бимодулей координатной алгебры и не создают отдельные копии
калибровочных бозонов. Поэтому допустима проверка архитектуры
\begin{equation}
 \mathcal A_{\mathrm{coord}}=\mathcal A_F^{\mathrm{SM}},
 \qquad
 \mathcal C_{\mathrm{fam}}
 \subseteq\pi(\mathcal A_F^{\mathrm{SM}})',
 \label{eq:v5-affine-commutant-route}
\end{equation}
где аффинная KO6-цепь является градуированным соответствием в коммутанте,
а не вторым координатным матричным множителем. Тогда стандарт-модельная
алгебра определяет физическую калибровочную группу, а семейный оператор
остаётся эквивариантным по всем 15 наблюдаемым компонентам.

Это пока направление, а не готовая спектральная тройка. Необходимо
проверить, определяет ли коммутантное соответствие самостоятельное
дифференциальное исчисление и действие Ходжа без возвращения внешнего
селектора. Литературная граница сверена с классификацией диаграмм
Краевского (arXiv:hep-th/9701081) и явным описанием поколений как
кратностей стандарт-модельного бимодуля (arXiv:1211.0825).

\begin{protocol}
\begin{enumerate}
 \item старое разложение $\mathbf1\oplus\mathbf3$:
 \textbf{восстановлено};
 \item единый бимодуль $(4\to3\to3)\times15\times2$:
 \textbf{построен};
 \item размерность 300:
 \textbf{получена без внешнего прямого сложения};
 \item KO6, нулевой и первый порядок:
 \textbf{выполнены};
 \item канонический физический угол с 45 частичными нулевыми состояниями:
 \textbf{получен};
 \item смешивания опоры с цепью:
 \textbf{строго стабилизированы};
 \item полная наблюдаемая координатная алгебра:
 \textbf{не построена};
 \item коммутантное размещение семейной цепи:
 \textbf{выделено для следующей проверки};
 \item единая кинетическая нормировка и повторный импульсный тест:
 \textbf{открыты};
 \item физическое замыкание:
 \textbf{не достигнуто}.
\end{enumerate}
\end{protocol}

Машинная проверка и сертификат сохранены в
\path{s2t_v5_affine_ko6_reference_corner_gate.py} и
\path{s2t_v5_affine_ko6_reference_corner_gate_results.json}.